\documentclass[11pt,a4paper]{article}
\usepackage[top=1.7cm,bottom=1.5cm,left=1.5cm,right=1.5cm,headsep=8pt,footskip=16pt]{geometry}
\usepackage{enumitem}
\usepackage{amssymb}
\usepackage{array}
\usepackage[table]{xcolor}
\usepackage{fancyhdr}
\usepackage{lastpage}
\usepackage{multicol}
\usepackage{polyglossia}
\setdefaultlanguage[calendar=gregorian, numerals=maghrib]{arabic}
\setotherlanguage{english}
\IfFontExistsTF{Amiri}{%
  \newfontfamily\arabicfont[Script=Arabic, Scale=1.08]{Amiri}
}{%
  \newfontfamily\arabicfont[Script=Arabic, Scale=1.08, Path=../fonts/, Extension=.ttf,
    UprightFont=*-Regular, BoldFont=*-Bold, ItalicFont=*-Italic,
    BoldItalicFont=*-BoldItalic]{Amiri}
}
\IfFontExistsTF{AlBattar}{%
  \newfontfamily\khatt[Script=Arabic, Scale=1.15]{AlBattar}
}{%
  \newfontfamily\khatt[Script=Arabic, Scale=1.15, Path=../fonts/,
    BoldFont=ae_AlBattar.ttf]{ae_AlBattar.ttf}
}
\renewcommand{\headrulewidth}{0.8pt}
\renewcommand{\footrulewidth}{0.8pt}
\fancyhf{}
\rhead{\khatt المدرسة العليا للأساتذة ببوزريعة}
\lhead{\khatt ملحقة بومرداس}
\lfoot{\khatt الأستاذ: د. بن عبد الرزاق يوسف}
\cfoot{\thepage/\pageref{LastPage}}
\pagestyle{fancy}
\parindent 0pt
\setlength{\columnsep}{0.6cm}
\setlength{\columnseprule}{0.4pt}
\setlist[enumerate,1]{label=\textbf{\arabic*.}, leftmargin=*, itemsep=5pt, topsep=2pt}
\begin{document}

\linespread{1.08}\selectfont
\begin{center}
{\LARGE\khatt\textbf{الاختبار السادس والعشرون في الذكاء الاصطناعي}}\\[4pt]
{\normalsize وضعيات إشكالية ــ أسئلة متعددة الاختيارات ــ الفصول من 1 إلى 6}
\end{center}
\vspace{-2pt}\rule{\textwidth}{0.6pt}\\[2pt]
\textbf{التعليمة:} يتضمن الاختبار خمس وضعيات إشكالية. اقرأ كل وضعية بعناية، ثم اختر التحليل أو التصرف الأنسب
من بين الاختيارات الخمسة؛ لكل وضعية إجابة صحيحة واحدة فقط.
اكتب \textbf{حرف} الإجابة الصحيحة في جدول الإجابات في آخر الاختبار.
لكل وضعية 4 نقاط (المجموع 20 نقطة).
\vspace{4pt}

\noindent
{\large\khatt\textbf{الوضعية 1}}\hfill{\small(الفصل 1)}\par\nopagebreak\vspace{3pt}
\setlength{\fboxsep}{7pt}\colorbox{gray!10}{\parbox{\dimexpr\linewidth-14pt}{ترجم طالب بأداة توليدية قصيدة إنجليزية قصيرة إلى العربية لنشرها في مجلة الكلية، فجاءت الترجمة سليمة المعنى، لكنها فقدت الإيقاع والجناس واللعب بالألفاظ الذي يقوم عليه النص الأصلي. ويرى زميله أن «الأداة ترجمت أحسن من أي مترجم». والمجلة تشترط ذكر اسم المترجم. ما التقدير الأدق؟}}\par\nopagebreak
\begin{itemize}[leftmargin=1.0cm, labelsep=6pt, itemsep=3pt, topsep=4pt, parsep=0pt]
  \item[\textbf{أ)}] يطلب الطالب من الأداة إضافة سجع في آخر كل سطر، ثم ينشر الترجمة باسمه دون مراجعة، لأن السجع يعوّض الإيقاع والجناس الأصليين تعويضًا كاملًا في كل الترجمات الشعرية.
  \item[\textbf{ب)}] الأداة تنقل المعنى غالبًا، لكنها قد تُسقط ما يقوم على خصائص اللغة الأصلية كالإيقاع والجناس؛ فتُعدّ ترجمتها مسودة يعيد الطالب صياغتها إبداعيًا، ويفصح عن دور الأداة عند النشر.
  \item[\textbf{ج)}] الترجمة فاشلة كليًا لأنها فقدت الإيقاع، والأدوات التوليدية لا تستطيع نقل المعنى أصلًا في النصوص الأدبية، فلا يمكن الاستفادة منها حتى مسودةً أولى يعيد الطالب صياغتها بنفسه.
  \item[\textbf{د)}] زميله محق، لأن سلامة المعنى هي كل ما تتطلبه ترجمة الشعر، والإيقاع والجناس زخارف لا تُترجم أصلًا، فتُنشر الترجمة كما هي باسم الطالب دون أي إشارة إلى الأداة المستعملة.
  \item[\textbf{هـ)}] يُنشر النص باسم الأداة وحدها مترجمًا، لأن الطالب لم يترجم شيئًا بنفسه، حتى لو أعاد صياغة الترجمة إبداعيًا، فكل نص مرّ بالأداة يُنسب إليها وحدها في كل الأحوال.
\end{itemize}
\par\vspace{6pt}
\noindent
{\large\khatt\textbf{الوضعية 2}}\hfill{\small(الفصل 2)}\par\nopagebreak\vspace{3pt}
\setlength{\fboxsep}{7pt}\colorbox{gray!10}{\parbox{\dimexpr\linewidth-14pt}{درّب فريق طلابي نموذجًا على بنك من ثلاثمئة سؤال إعراب مع إجاباتها، وجرّبوه على الأسئلة نفسها فأجاب عنها كلها بدقة. لكنه عند تغيير كلمة واحدة في الجملة أعاد الإعراب القديم نفسه أحيانًا. ويستعد الفريق لعرض النموذج على أساتذة القسم. ما التقييم الأدق لهذه النتيجة؟}}\par\nopagebreak
\begin{itemize}[leftmargin=1.0cm, labelsep=6pt, itemsep=3pt, topsep=4pt, parsep=0pt]
  \item[\textbf{أ)}] النموذج ممتاز، لأن دقته الكاملة على الأسئلة نفسها دليل على أنه أتقن الإعراب، وتغيير كلمة واحدة لا يغيّر الإعراب في العربية عادة، فلا حاجة إلى اختبار إضافي قبل العرض.
  \item[\textbf{ب)}] النموذج حفظ الأسئلة، والحل أن يُطلب منه في كل مرة «لا تعتمد على الحفظ»، لأن هذه التعليمة تجعله يطبّق القواعد على الجمل الجديدة دون حاجة إلى توسيع البيانات أو تنويعها.
  \item[\textbf{ج)}] النموذج حفظ الأسئلة، والحل تقليص البنك إلى مئة سؤال، لأن البيانات القليلة تجبر النموذج على التعميم دائمًا، ثم يُعاد اختباره على الأسئلة المئة نفسها التي تدرّب عليها.
  \item[\textbf{د)}] الخلل في تغيير الكلمات لا في النموذج، فيُمنع المستخدمون من تعديل جمل البنك، ويُعرض النموذج على الأساتذة بوصفه أداة إعراب عامة، لأن دقته على البنك مثبتة تمامًا.
  \item[\textbf{هـ)}] النموذج حفظ الأسئلة وإجاباتها أكثر مما تعلّم قواعد قابلة للتعميم، والتجربة على بيانات التدريب نفسها لا تكشف ذلك؛ فيُختبر على جمل جديدة لم يرها، وتُوسَّع البيانات وتُنوَّع.
\end{itemize}
\par\vspace{6pt}
\noindent
{\large\khatt\textbf{الوضعية 3}}\hfill{\small(الفصل 4)}\par\nopagebreak\vspace{3pt}
\setlength{\fboxsep}{7pt}\colorbox{gray!10}{\parbox{\dimexpr\linewidth-14pt}{رفع أستاذ الكتاب المدرسي كاملًا إلى أداة ذكاء اصطناعي، وطلب ملخصات لكل دروسه، ثم جمعها في كتيّب أراد بيعه للتلاميذ وأوليائهم في المكتبات. ويرى أن الملخصات «من إنتاج الأداة»، فلا علاقة لها بحقوق مؤلفي الكتاب المدرسي. ويستعد لطبع ألف نسخة قبل بداية الفصل الثاني. ما التقدير الأدق؟}}\par\nopagebreak
\begin{itemize}[leftmargin=1.0cm, labelsep=6pt, itemsep=3pt, topsep=4pt, parsep=0pt]
  \item[\textbf{أ)}] الملخصات مشتقة من مضمون كتاب محمي، وبيعها دون إذن أصحاب الحقوق يثير إشكالًا قانونيًا وأخلاقيًا مهما كانت الأداة المستعملة؛ فيستأذن الجهة المعنية أو يكتفي باستعمال تعليمي غير تجاري في حدود المسموح.
  \item[\textbf{ب)}] يكفي أن يكتب على الغلاف «لُخّص بالذكاء الاصطناعي» ليصبح البيع مشروعًا، لأن الإفصاح عن الأداة يحوّل الحقوق إليها، ولا يحتاج بعد ذلك إلى استئذان أي جهة.
  \item[\textbf{ج)}] يجوز البيع إذا خفّض السعر إلى أدنى حد، لأن حقوق المؤلف تتعلق بالربح الكبير فقط، والبيع بسعر رمزي لا يُعدّ استغلالًا تجاريًا للكتاب المدرسي مهما كان عدد النسخ، فيبقى الكتيّب في متناول الجميع.
  \item[\textbf{د)}] لا إشكال، لأن الأداة هي التي لخّصت، والتلخيص الآلي يُنتج نصًا جديدًا لا علاقة له بحقوق المؤلفين، فيجوز بيع الكتيّب باسم الأستاذ في كل المكتبات دون أي إذن، ويحق له أن يحتفظ بكل الأرباح.
  \item[\textbf{هـ)}] الإشكال الوحيد في رفع الكتاب إلى الأداة، أما البيع فمشروع، لأن الملخصات أقصر من الكتاب، والنص القصير المشتق من كتاب محمي لا تشمله حقوق المؤلف في أي حال.
\end{itemize}
\par\vspace{6pt}
\noindent
{\large\khatt\textbf{الوضعية 4}}\hfill{\small(الفصل 5)}\par\nopagebreak\vspace{3pt}
\setlength{\fboxsep}{7pt}\colorbox{gray!10}{\parbox{\dimexpr\linewidth-14pt}{انتشر على مواقع التواصل منشور يقول إن «العربية تضم اثني عشر مليون كلمة، وهي أغنى لغات العالم بفارق كبير». سأل طالب أداة ذكاء اصطناعي فأكدت الرقم وأضافت مقارنات بلغات أخرى. ويريد إدراج المعلومة في عرضه أمام زملائه. وأستاذه يشترط توثيق كل رقم يُعرض. ما المنهج الأنسب للتحقق منها؟}}\par\nopagebreak
\begin{itemize}[leftmargin=1.0cm, labelsep=6pt, itemsep=3pt, topsep=4pt, parsep=0pt]
  \item[\textbf{أ)}] يعتمد الرقم لأن الأداة أكدته وأضافت مقارنات تفصيلية، والتفاصيل الإضافية دليل على أن الأداة رجعت إلى دراسات لسانية دقيقة قبل الإجابة، فلا حاجة إلى البحث عن أصل الرقم.
  \item[\textbf{ب)}] يرفض المعلومة فورًا دون أي بحث، لأن كل ما يُنشر على مواقع التواصل خاطئ بالضرورة، ولا يمكن أن يتضمن أي جزء صحيح، حتى لو وجد له أصلًا في مصادر محكّمة.
  \item[\textbf{ج)}] يقرأ قراءة جانبية: يبحث عن أصل الرقم وطريقة حسابه، وهل يعدّ الصيغ المشتقة كلمات مستقلة، ويقارن بمصادر لسانية محكّمة، لأن تأكيد الأداة قد يعكس انتشار الادعاء لا صحته.
  \item[\textbf{د)}] يسأل أداة ثانية، فإن أكدت الرقم أيضًا اعتمده ونسبه إلى «الدراسات اللسانية»، لأن اتفاق أداتين يغني عن تحديد الدراسة الأصلية وطريقة حساب الكلمات فيها.
  \item[\textbf{هـ)}] يعتمد الرقم لأنه منتشر في منشورات كثيرة، وكثرة التداول على مواقع التواصل دليل على صحة المعلومة، خاصة إذا أكدتها أداة ذكاء اصطناعي في الوقت نفسه دون تردد.
\end{itemize}
\par\vspace{6pt}
\noindent
{\large\khatt\textbf{الوضعية 5}}\hfill{\small(الفصل 6)}\par\nopagebreak\vspace{3pt}
\setlength{\fboxsep}{7pt}\colorbox{gray!10}{\parbox{\dimexpr\linewidth-14pt}{تعدّ أستاذة بطاقات مصوّرة لمفاهيم قيمية: «الأمانة، العدل، الكرم، الصبر» لتلاميذ السنة الثانية ابتدائي، باستعمال أداة توليد صور. جاءت الصور جميلة، لكن «الكرم» صُوّر بمائدة فاخرة، و«العدل» بميزان فقط، و«الصبر» بشخص حزين. وتريد بطاقات تبني فهمًا صحيحًا للقيم. وستعلّق البطاقات في القسم طوال السنة. ما التصرف الأنسب؟}}\par\nopagebreak
\begin{itemize}[leftmargin=1.0cm, labelsep=6pt, itemsep=3pt, topsep=4pt, parsep=0pt]
  \item[\textbf{أ)}] تطلب صورًا أكثر عددًا لكل قيمة وتختار أجملها دون مراجعة دلالتها، لأن الجمال البصري هو ما يجذب الأطفال، والدلالة القيمية للصورة تتحقق تلقائيًا عند كتابة اسم القيمة تحتها، فتنجذب إليها أنظار التلاميذ أكثر.
  \item[\textbf{ب)}] تكتب اسم القيمة فقط في الطلب بخط كبير، لأن الأداة تفهم المفاهيم المجردة كالإنسان، وتختار تلقائيًا أنسب تمثيل لها في عمر التلاميذ، دون حاجة إلى وصف موقف أو مراجعة، ويتعلمونها دون جهد في التأويل.
  \item[\textbf{ج)}] تعتمد الصور كما هي، لأن الرموز المشهورة كالميزان والمائدة هي أسهل طريق لتعليم القيم، والأطفال يفهمون الرموز المجردة بسهولة، فلا حاجة إلى مواقف ملموسة أو نقاش.
  \item[\textbf{د)}] تصف في الطلب مواقف ملموسة من حياة الطفل تجسّد كل قيمة، كتلميذ يعيد محفظة ضائعة، وتراجع الصور مع زميلة لتتأكد من أنها لا تختزل القيمة أو تشوّهها، وتصاحبها بقصة قصيرة أو نقاش.
  \item[\textbf{هـ)}] تتخلى عن الصور وتكتب تعريفات معجمية للقيم على البطاقات، لأن الصور لا يمكن أن تمثل المفاهيم المجردة أبدًا، والتعريفات المعجمية أنسب لتلاميذ السنة الثانية ابتدائي.
\end{itemize}
\par\vspace{6pt}
\begin{center}
{\khatt\textbf{جدول الإجابات}}\\[3pt]
\renewcommand{\arraystretch}{1.6}
\begin{tabular}{|>{\centering\arraybackslash}m{1.9cm}|*{5}{>{\centering\arraybackslash}m{1.2cm}|}}
\hline
\textbf{رقم السؤال} & \textbf{1} & \textbf{2} & \textbf{3} & \textbf{4} & \textbf{5} \\ \hline
\textbf{الإجابة} &  &  &  &  &  \\[8pt] \hline
\end{tabular}

\vspace{8pt}
\renewcommand{\arraystretch}{2.2}
\begin{tabular}{|>{\centering\arraybackslash}m{7.5cm}|>{\centering\arraybackslash}m{3.2cm}|>{\centering\arraybackslash}m{3.2cm}|}
\hline
\textbf{الاسم واللقب} & \textbf{الفوج} & \textbf{العلامة} \\ \hline
 & & \hfill /20 \\ \hline
\end{tabular}
\end{center}
\end{document}
